\Needspace{10\baselineskip}
\section{酉算子的谱分解}\label{sec:yosida-11-04}
\noindent\textit{原书起始页：中译本 319；英文版 306。}

\begin{theorem}[单位圆上的谱积分]\label{thm:y11-unitary-pvm}
若 $U$ 酉，则其谱测度支于单位圆 $\mathbb T$，且对所有 $n\in\mathbb Z$，
\[
 U^n=\int_{\mathbb T}z^n\,\dd E(z),\qquad U^*=\int_{\mathbb T}\bar z\,\dd E(z).
\]
反之，单位圆上总质量为 $I$ 的投影值测度的坐标积分是酉算子。
\end{theorem}
\begin{proof}
$U$ 正规，$U^*U=I$ 使连续演算中 $|z|^2=1$ 于谱上；也可对 $|\lambda|>1$ 用 $U$ 的 Neumann 级数，对 $|\lambda|<1$ 用 $U^{-1}$ 的级数得谱包含。各幂由演算乘法得到，反向由积分 $z\bar z=1$ 得 $U^*U=UU^*=I$。
\end{proof}

选定 Borel 辐角 $\operatorname{Arg}:\mathbb T\to[0,2\pi)$，定义 $B=\operatorname{Arg}(U)$。则 $B$ 为有界自伴算子，$U=e^{iB}$，但这个对数依赖所选分支，不能据此宣称唯一。若用实参数谱族 $F_t=E(\{z:\operatorname{Arg}z\le t\})$，应把 $z=1$ 的质量计入 $t=0$；原书使用另一端点记法时须整体换算，不能重复计算该原子。

\begin{example}[平移的谱表示]
在 $\ell^2(\mathbb Z)$ 上取双边移位 $(Sx)_n=x_{n-1}$。圆周取归一化测度 $\dd\theta/(2\pi)$，则酉 Fourier 级数映射 $\mathcal Fx(\theta)=\sum_nx_ne^{in\theta}$ 将它变为乘 $e^{i\theta}$。因此谱为 $\mathbb T$，谱投影为 $\mathcal F^{-1}1_B(e^{i\theta})\mathcal F$。单点测度为零，所以没有特征值。
\end{example}

\begin{exercise}[谱投影的原子]\label{ex:y11-unitary-atom}
证明 $E(\{z_0\})H=\ker(U-z_0I)$。
\end{exercise}
\begin{solution}
由 $\|(U-z_0)x\|^2=\int|z-z_0|^2\,\dd\mu_x$，该范数为零当且仅当 $\mu_x$ 集中于 $z_0$。这又等价于 $E(\{z_0\})x=x$。
\end{solution}

\begin{exercise}[酉对数不唯一]\label{ex:y11-unitary-log}
给出两个不同的有界自伴算子 $B_1,B_2$，使 $e^{iB_1}=e^{iB_2}=I$。
\end{exercise}
\begin{solution}
取 $B_1=0$，$B_2=2\pi P$，其中 $P$ 为任意非零正交投影。连续演算给出 $e^{2\pi iP}=(I-P)+P=I$。
\end{solution}
